\section{Parameters}
\label{sec:parameters}

Two sections below \ui{REGION OF INTEREST} hold the remaining settings of a run
(Figure~\ref{fig:parameters}). \ui{PARAMETERS},
open by default, sets how the network runs; \ui{ADVANCED}, closed by default (click its title to
open it), sets how each image pair is cut into windows for the network. Every new session starts
from the defaults in Table~\ref{tab:parameters}; a project file keeps the values you set
(Section~\ref{sec:projects}).
% CHECK: after File > Open Project the ADVANCED controls and the Interval box keep their previous values on screen, although the run uses the project's values (AdvancedPanel has no sync from AppState); reproduced offscreen

\begin{table}[htbp]
  \centering
  \caption{The controls of \ui{PARAMETERS} and \ui{ADVANCED}.}
  \label{tab:parameters}
  \begin{tabularx}{\textwidth}{>{\raggedright\arraybackslash}p{2.6cm}
                               >{\raggedright\arraybackslash}p{5.2cm}
                               >{\raggedright\arraybackslash}X}
    \toprule
    Control & Default & Range or choices \\
    \midrule
    \ui{Iterations}   & 20 iterations & 1 to 100 iterations \\
    \ui{Smoothing}    & \ui{Off} (0 px) & 0 to 50 px, in steps of 0.5 px \\
    \ui{Device}       & \ui{Automatic} & \ui{Automatic}, \ui{GPU (CUDA)}, \ui{CPU} \\
    \midrule
    \ui{Margin}       & \ui{Use the Model Reach} ticked ($32\,s$ px) & 0 to 4096 px \\
    \ui{GPU Memory}   & 0.85 (fraction of the free memory) & 0.05 to 1.00, in steps of 0.05 \\
    \ui{Tile Blending} & 32 px & 0 to 512 px \\
    \ui{Correlation}  & \ui{Automatic} & \ui{Automatic}, \ui{On Demand}, \ui{Stored Volume} \\
    \bottomrule
  \end{tabularx}
\end{table}

\screenshot[0.45\textwidth]{fig-parameters.png}{The \ui{PARAMETERS} and \ui{ADVANCED} sections,
both open, with their default values.}{fig:parameters}

Hovering over a label shows what it does. A spin box or a list reacts to the mouse wheel only
after you click it, so scrolling the sidebar does not change a value by accident.

\subsection{PARAMETERS}

\begin{description}[leftmargin=1.2cm,style=nextline]
  \item[Iterations]
    The number of refinement iterations of the network (default 20). The network refines its
    estimate of the displacement once per iteration, so the computing time grows with the
    number. \ui{Suggest a Model} uses the same value, and the card on the canvas shows it under
    \ui{ITERATIONS}. Keep the default unless you have a reason to change it; if you lower it to
    shorten a run, compare the residual (Section~\ref{sec:results}) with that of a run at the
    default before you rely on the result.
  \item[Smoothing]
    Gaussian smoothing of the tracked displacement, set as its standard deviation $\sigma$ in
    pixels. It is \ui{Off} by default, so the results hold the displacement as the network
    measures it. Smoothing coarsens the spatial resolution: features of the field smaller than
    a few $\sigma$ are flattened. The smoothing stays inside the region of interest: empty points
    and points outside the region are not used and stay empty, and a point with almost no valid
    neighbours (less than 5\,\% of the weight of the Gaussian) becomes empty. The smoothed
    displacement is what the results, the strain window and the exports receive; in
    incremental tracking it is also what later frames are composed with. For strain alone,
    prefer \ui{Pre-fit smoothing} in the strain window (Section~\ref{sec:strain}), which feeds
    only the strain fit and leaves the displacement as it is.
  \item[Device]
    Where the network runs. \ui{Automatic} uses the GPU when PyTorch sees one (an NVIDIA GPU
    with CUDA) and the CPU otherwise; on macOS RAFTcorr runs on the CPU. \ui{GPU (CUDA)} on a
    computer without CUDA stops the run with an error. Choose \ui{CPU} to leave the GPU to
    another program. \menu{Help}{Self-test} checks the installed models on the device you choose
    (Section~\ref{sec:installing-check}).
\end{description}

\subsection{ADVANCED}

For every pair, the network receives a \emph{window} of both images: the bounding box of the
region of interest of the reference frame, grown by a margin on every side. When that window
needs more memory than the device offers, RAFTcorr cuts the region into \emph{tiles}. The plan is
automatic; Section~\ref{sec:tiling} describes it, its line in the \ui{LOG} and what happens when the
memory runs out. The four controls below adjust it.

\begin{description}[leftmargin=1.2cm,style=nextline]
  \item[Margin and Use the Model Reach]
    The margin is the image context added around the bounding box, so that a point that moves
    stays inside the window the network sees. It must cover the largest displacement of the
    region of interest. With \ui{Use the Model Reach} ticked (default), the margin is the reach
    of the selected model: 32, 64, 128 or 256~px for RAFTcorr-s1, -s2, -s4 or -s8; the box is
    greyed out and its value is not used. Untick it to type a margin from 0 to 4096~px. A
    smaller margin than the reach saves memory when you use a model with a large $s$ on small
    motion; a margin smaller than the displacement loses the points that move out of the
    window. The window never extends past the image, and when tiles are needed the margin can
    be reduced to fit the memory, in which case the tiling line of the log says so.
  \item[GPU Memory]
    The fraction of the free GPU memory one window may use, measured before each pair (default
    0.85; RAFTcorr also keeps 384~MB free). A larger region is cut into tiles to fit. Lower it
    when another program needs the GPU or when runs fail for lack of memory; raise it when the
    GPU is free and you want fewer tiles. On the CPU it has no effect: windows are planned for a
    fixed budget of 6~GB.
  \item[Tile Blending]
    The width of the overlap between neighbouring tiles, where their results are blended with
    linear weights (default 32~px). It matters only when the region is cut into tiles. A wider
    overlap spreads the change from one tile to the next over more pixels and needs more tiles;
    0~px lets the tiles meet edge to edge. The default rarely needs a change.
  \item[Correlation]
    How the network computes its correlation lookup. Both ways give the same displacement; they
    differ in memory and speed. \ui{Stored Volume} computes the correlation between all pairs of
    feature points of the window once and keeps it, so its memory grows with the square of the
    window area; it is the faster choice where it fits. \ui{On Demand} computes each lookup when it is needed,
    with memory in proportion to the window area. \ui{Automatic} uses \ui{Stored Volume} for
    RAFTcorr-s8 and \ui{On Demand} for -s1, -s2 and -s4. Leave it at \ui{Automatic}; with
    RAFTcorr-s8 and little memory, \ui{On Demand} allows larger windows and fewer tiles.
\end{description}
